%% =====================================================================
%% Manuscript for Aerospace Science and Technology (Elsevier)
%% Uses the official Elsevier elsarticle class (elsarticle-num style).
%% Results in Section 8 are recomputed from the plant-truth CSV logs under
%% Tools/tilt_hexa_30kg/results/MPC (see paper/README.md). All airframe
%% numbers in Section 3 are REFERENCE_SEED_NOT_MEASURED for the 30 kg design.
%% Author/affiliation fields are placeholders to be completed before submission.
%% =====================================================================
\documentclass[preprint,3p,11pt]{elsarticle}

\usepackage{amsmath,amssymb,amsthm}
\usepackage{graphicx}
\usepackage{booktabs}
\usepackage{multirow}
\usepackage{array}
\usepackage{url}
\usepackage{hyperref}
\usepackage{xcolor}
\usepackage{enumitem}

% Lightweight siunitx-compatible macros (avoids a missing siunitx package).
\newcommand{\SI}[2]{\ensuremath{#1\,#2}}
\newcommand{\si}[1]{\ensuremath{#1}}
\newcommand{\kilogram}{\mathrm{kg}}
\newcommand{\newton}{\mathrm{N}}
\newcommand{\meter}{\mathrm{m}}
\newcommand{\kilo}{\mathrm{k}}
\newcommand{\watt}{\mathrm{W}}
\newcommand{\second}{\mathrm{s}}
\newcommand{\pascal}{\mathrm{Pa}}
\newcommand{\per}{/}
\newcommand{\squared}{^{2}}
\newcommand{\degree}{^{\circ}}

\hypersetup{colorlinks=true,linkcolor=black,citecolor=blue,urlcolor=blue}

\journal{Aerospace Science and Technology}

\begin{document}

\begin{frontmatter}

\title{Corridor-aware conversion scheduling and unified model predictive control for a fully-tilting hexacopter: an open-source SITL companion benchmark in the ArduPilot framework}

\author[a]{Author~One}
\author[a]{Author~Two}
\author[a,b]{Corresponding~Author\corref{cor1}}
\cortext[cor1]{Corresponding author.}
\affiliation[a]{organization={School of Automation, Nanjing University},
  addressline={Nanjing, China}}
\affiliation[b]{organization={(Corresponding affiliation)}}

\begin{abstract}
Fully-tilting hexacopters combine the hover capability of a multirotor with
the efficient cruise of a fixed-wing aircraft, but the continuous transfer of
lift between rotors and wing during conversion makes their full-envelope
control difficult. Production autopilots, including the stock ArduPilot
QuadPlane, address this problem with a switched multirotor/fixed-wing
architecture and an open-loop, fixed-rate tilt schedule, which can induce
altitude loss, attitude transients, and a large tuning burden. This paper does
not claim to invent unified control; rather, it targets three issues that
remain open for the \emph{fully-tilting hexacopter}, a strongly over-actuated
configuration with six independently tilting rotors. First, a time-varying
control-allocation layer is developed around the \emph{attainable wrench set}
(AWS), whose inscribed-sphere radius is evaluated online along the
transition to detect saturation and redistribute the six throttle commands,
six tilt angles, and the aerodynamic control surfaces. Second, a
corridor-constrained optimal control problem (OCP) generates the
forward and backward conversion reference, and a single unified model
predictive controller (MPC) tracks it across the whole envelope with zero mode
switches; an AWS-margin feasibility signal is defined to trigger conservative
re-planning if the manoeuvre becomes infeasible, although this path was not
entered in the tested envelope because the constrained allocator kept a
positive margin. A bounded, model-free position-error
integral rejects constant plant mismatch without the instability of a
model-based disturbance observer. Third, the complete pipeline is released
as an open, reproducible software-in-the-loop (SITL) companion benchmark built
in the ArduPilot framework for a \SI{30}{\kilogram} vehicle, and is compared
under a standardized scenario matrix against an INDI--WLS baseline. In
plant-truth simulation the controller completes the full hover--cruise--hover
mission with zero mode switches, peak attitude below $18^\circ$,
forward/backward conversions of 14.0/20.1\,s and 21.0/26.4\,kJ of rotor
induced energy, and passes all nine truth checks in eight fixed robustness
scenarios (gust, steady crosswind, $+15\%$ mass, CG offset,
thrust/surface/inertia mismatch) and a 20-run Monte-Carlo campaign. A tuned
INDI--WLS baseline tracks more tightly in nominal conditions; the proposed
controller does not claim tracking superiority, and its demonstrated value is
the zero-switch unified structure, explicit constraint handling, and the
feasibility guarantee of the corridor. Removing the corridor destabilizes the
backward leg, raising its altitude deviation from $0.98$ to $7.0$~m and its
energy from $26.4$ to $45.6$~kJ and failing the truth checks. The inner QP runs
at a mean of $0.5$~ms and a 99th percentile below $2.4$~ms against a $30$~ms
period, with the C++ allocator at $49$--$74\,\mu$s. The benchmark, parameters,
and scenario files are provided to make the study reproducible and to serve as
a common baseline for future convertible-aircraft controllers.
\end{abstract}

\begin{keyword}
Tilt-rotor hexacopter \sep Model predictive control \sep Control allocation
\sep Attainable wrench set \sep Conversion corridor \sep ArduPilot SITL
\end{keyword}

\end{frontmatter}

% =====================================================================
\section{Introduction}\label{sec:intro}

Convertible vertical take-off and landing (VTOL) unmanned aerial vehicles
(UAVs) extend the endurance of multirotors by adding a lifting wing, and
improve the versatility of fixed-wing aircraft by permitting hover. Among the
many convertible layouts reviewed by Ducard and Allenspach~\cite{ducard2021review},
the \emph{tilt-rotor} family rotates the propulsion units between a vertical
(hover) and a horizontal (cruise) orientation. This paper studies a
\emph{fully-tilting hexacopter}: six rotors are each mounted on an
independently actuated nacelle, rotating about the spanwise axis from
\SI{0}{\degree} (thrust vertical) to \SI{90}{\degree} (thrust horizontal),
while a fixed wing and conventional tail surfaces provide lift and control in
cruise. The vehicle considered has a maximum take-off mass of
\SI{30}{\kilogram}.

The principal difficulty of such aircraft is the conversion corridor. During
forward transition the rotors must tilt forward to gain airspeed while the wing
progressively carries the weight; during backward transition the wing unloads
while the rotors must re-establish vertical lift. Neither extreme is
problematic, but the intermediate region combines rapidly changing control
authority, propeller-wing aerodynamic interaction, and near-saturation
actuators. The autopilot must simultaneously avoid the low-speed stall
boundary and the high-speed power/load boundary, while tracking a feasible
attitude and altitude profile.

The architecture used by the open-source ArduPilot autopilot is
representative of production practice. In its QuadPlane framework, a
multirotor controller and a fixed-wing controller are switched in and out, and
the tilt servos are driven through an open-loop schedule with a fixed tilt
rate~\cite{ardupilot2024tilt,ardupilot2024transitions}. This is robust and
well tested, but it has three structural drawbacks: (i) the controller is
discontinuous across the envelope, which produces transients at the switching
points; (ii) the tilt schedule is fixed a priori and is not adapted to wind,
mass, or actuator saturation; and (iii) all tilt servos are treated as a
single synchronized logical tilt, which does not exploit the independent
per-rotor tilting authority of a fully-tilting hexacopter.

A large body of academic work has already addressed unified, full-envelope
control. Bauersfeld \emph{et al.}~\cite{bauersfeld2021mpc} demonstrated, in
outdoor flight, a model predictive controller that accepts a pilot velocity
command and directly computes rotor speeds and tilt angles across the entire
envelope. Allenspach and Ducard~\cite{allenspach2021nmpc} combined a unified
nonlinear MPC (NMPC) with multi-stage control allocation on a real
propeller-tilting UAV; Rohr \emph{et al.}~\cite{rohr2021nmpc} used NMPC for
rapid tilt-wing transitions; Mancinelli \emph{et al.}~\cite{mancinelli2025unified}
validated a unified incremental nonlinear controller with angle-of-attack
protection on a dual-axis tilting quad-plane; and Chen \emph{et
al.}~\cite{chen2024unified} reported a switching-free unified MPC. Incremental
nonlinear dynamic inversion (INDI) offers a closely related, model-robust
alternative~\cite{smeur2020incremental,smeur2018cascaded,smeur2016adaptive,wang2019stability},
and has been applied to tailsitters and tilt-rotors. Consequently, the
\emph{concept} of unified control is not novel and is not claimed as a
contribution here.

Three gaps nevertheless remain. First, almost all unified controllers target
quad-scale tilt-rotors, tilt-wings, or tailsitters; the fully-tilting
hexacopter, six rotors and six tilt servos plus four control surfaces (two
ailerons and two V-tail ruddervators), i.e. sixteen actuators for six degrees
of freedom, has a strongly time-varying,
heavily redundant control-allocation problem that has mostly been studied for
hover fault tolerance rather than across the full envelope. Second, optimal
transition scheduling within a conversion corridor has been developed for
manned tilt-rotors and tilt-wings, but is rarely coupled, in closed loop, with
the real-time feasible-wrench information of the tracking controller. Third,
although software-in-the-loop (SITL) augmentation and benchmarking of open
autopilots has been shown for fixed-wing and conventional multirotors, no
open, reproducible SITL benchmark currently covers a fully-tilting hexacopter
with a unified controller and a corridor-optimal transition.

This paper therefore makes three targeted contributions.

\begin{enumerate}[label=(\roman*),leftmargin=2.2em]
  \item \textbf{Time-varying AWS allocation.} A control-allocation layer is
  formulated for the fully-tilting hexacopter in which the control
  effectiveness matrix is updated online with tilt angle and airspeed, and a
  constrained QP allocates the virtual wrench to the sixteen actuators with
  position and rate limits as hard constraints. The attainable wrench set
  (AWS) and its inscribed-sphere radius are evaluated along the transition to
  provide a scalar feasibility margin, and the architecture specifies a
  wrench-hedging action when that margin vanishes.
  \item \textbf{Corridor-feasible reference and unified MPC.} A
  corridor-constrained trim optimization (computed offline) supplies a
  feasible-by-construction forward and backward conversion reference, and a
  single unified model predictive controller tracks it across the whole
  envelope with zero mode switches. A feasibility signal is defined to couple
  the two layers and trigger conservative re-planning; this re-planning path is
  designed and specified but is not triggered in the tested envelope, and the
  demonstrated online feasibility comes from the constrained QP and the
  corridor reference.
  \item \textbf{Open SITL companion benchmark.} The method is implemented as a
  companion controller around a high-rate nonlinear plant in the ArduPilot
  framework for a \SI{30}{\kilogram} fully-tilting hexacopter and released,
  with parameters, scenarios, and an INDI--WLS engineering baseline, as a
  reproducible benchmark. The closed-loop quantitative results are generated
  by the companion plant; the native ArduPilot \texttt{arduplane} SITL target
  and firmware hook are released as engineering artifacts and are identified
  as the next integration step rather than as the source of the reported
  numbers. A single-stick velocity command is included as a design property,
  with the number of mode switches as its evaluated metric.
\end{enumerate}

The remainder of the paper is organized as follows. Section~2 reviews related
work. Section~3 describes the platform and the two models used for control
and simulation. Section~4 develops the time-varying AWS allocation, Section~5
the corridor OCP, and Section~6 the unified MPC and the two-layer
integration. Section~7 details the ArduPilot SITL implementation, Section~8
defines the simulation campaign, and Section~9 discusses limitations.
Section~10 concludes.

% =====================================================================
\section{Related work}\label{sec:related}

\subsection{Unified full-envelope control of convertible aircraft}
Two controller families dominate. Optimization-based methods compute
attitude, rotor speeds, and tilt angles in a single predictive
problem~\cite{bauersfeld2021mpc,allenspach2021nmpc,rohr2021nmpc,chen2024unified}.
Bauersfeld \emph{et al.}~\cite{bauersfeld2021mpc} is the most direct prior
work: a pilot velocity command is mapped to rotor thrusts and tilt angles by
an MPC over the entire envelope, with outdoor flight results and an explicit
comparison against a binary-switch strategy. Sensor-based incremental methods
(INDI) instead linearize about the current state and use angular acceleration
feedback, which makes them robust to aerodynamic uncertainty; they have been
validated on tailsitters, tilt-rotors, and dual-axis
tilters~\cite{smeur2020incremental,smeur2018cascaded,smeur2016adaptive,wang2019stability,mancinelli2025unified,bhardwaj2018integrated,lovell2023tailsitter,sun2021incremental}.
Other recent approaches include iterative-LQR MPC for tilt-propulsion
eVTOL~\cite{xia2022ilqr}, robust allocation for a quad-plane
RAMC~\cite{cardoso2020ramc}, generic NMPC for multirotors with fixed and tilted
propellers~\cite{bicego2020nmpc}, unified attitude control in
ACC~\cite{belak2023unified}, MPC for a tilt-rotor tricopter~\cite{prach2018mpc},
neural-network approximations of MPC~\cite{ducard2024neural,jiang2022nnmpc},
dynamic-inversion transition
controllers~\cite{saetti2025dynamic,jeong2025longitudinal,li2026asynchronous},
and composite tilt-rotor modelling~\cite{liang2024composite}. Feasible control
allocation has been coupled with NMPC for a tiltrotor
quadrotor~\cite{shayan2024nmpc}, the closest predictive-allocation precedent;
that work targets a four-rotor machine with a fixed envelope and neither
constructs a conversion corridor nor closes a feasibility loop to an outer
planner. The present work uses MPC because it handles actuator constraints and
the conversion corridor naturally; the novelty is not the predictive control
itself but its coupling with the real-time AWS and its realization on a
fully-tilting hexacopter with a corridor-feasible reference.

\subsection{Transition scheduling and the conversion corridor}
The conversion corridor delimits the safe (airspeed, tilt) combinations.
Yu \emph{et al.}~\cite{yu2022pilot} optimized the XV-15 conversion to
alleviate pilot workload, explicitly including corridor constraints, power,
attitude, and a multi-objective cost. Kang \emph{et al.}~\cite{kang2024minimum}
computed minimum-energy conversion within a predefined corridor; May
\emph{et al.}~\cite{may2025backward} studied backward transition under
failure; and Milz \emph{et al.} examined unified pilot control and simplified
vehicle operations~\cite{milz2026taming} and robust NDI for a tandem
tilt-wing~\cite{milz2024tandem}. Tilt-wing cargo optimization was treated
in~\cite{chen2020tiltwing}; forward and backward transition trajectory
optimization with constant-altitude and stall constraints was studied
in~\cite{panish2024tiltwing}; and a multi-objective constrained tilt corridor
with a dynamic-imbalance and power boundary was proposed for a compound
tilt-rotor in~\cite{zhuang2025mctc}. These studies establish corridor-optimal
scheduling, but generally treat the corridor as a fixed offline envelope and
do not feed the tracking controller's instantaneous feasibility back to the
planner.

\subsection{Control allocation for over-actuated tilt multirotors}
Control allocation maps a desired virtual wrench to redundant actuators under
saturation; the field is reviewed in~\cite{johansen2013survey,blaha2023survey},
with adaptive constrained variants in~\cite{tohidi2020adaptive}. Tilted and
tilting hexacopters have been studied mainly for hover fault tolerance and
full actuation: Giribet \emph{et al.}~\cite{giribet2016tilted} designed a
tilted-rotor hexacopter for fault tolerance; Pose \emph{et al.}~\cite{pose2017hexacopter}
analysed hexacopter allocation after rotor loss; fast allocation for a heavily
over-actuated quad tilt-rotor was developed in~\cite{santos2021fast}; and
sliding-mode, self-scheduled, robust, and non-negative-least-squares allocators
were reported
in~\cite{nguyen2022smc,nguyen2018selfscheduled,baldini2023ftc,zheng2026safety}.
Yang \emph{et al.}~\cite{yang2026biaxial} used the inscribed-sphere metric of
the AWS to preserve full six-degree-of-freedom actuation of a biaxial-tilt
hexacopter under passive fault tolerance; Mancinelli
\emph{et al.}~\cite{mancinelli2025ftc} addressed over-actuated hybrid VTOL
fault tolerance; and controllability after a single rotor failure was analysed
in~\cite{du2015controllability}. The AWS metric is therefore established, but
it has been used as an offline or hover-only diagnostic. We evaluate it
\emph{online along the full transition} and close the loop to the planner.

\subsection{Open-source autopilots and SITL validation}
Open flight stacks are surveyed in~\cite{aliane2024opensource}. ArduPilot has
been used as a platform for adaptive augmentation with SITL
validation~\cite{baldi2022ardupilot}, adaptive formation
flight~\cite{yang2019formation}, and digital-twin
modelling~\cite{valencia2025digitaltwin}; the PX4 stack was used for fractional
control of an omnidirectional hexarotor~\cite{deoca2025fractional}; HIL and
VTOL swarm simulation platforms were presented
in~\cite{dai2021rflysim,sethi2022swarm}. None of these covers a fully-tilting
hexacopter with a unified, corridor-aware controller. Our benchmark fills this
gap and is deliberately built on the widely deployed ArduPilot SITL.

\subsection{Positioning relative to the closest prior work}
Relative to Bauersfeld \emph{et al.}~\cite{bauersfeld2021mpc}, the differences
are: (i) the platform is a fully-tilting hexacopter (sixteen actuators) rather
than a quad-scale tilt-rotor; (ii) the allocation is made explicitly
time-varying and feasibility-aware through the AWS; and (iii) the result is an
open ArduPilot SITL benchmark rather than a proprietary flight stack. Relative
to Yu \emph{et al.}~\cite{yu2022pilot}, the corridor OCP is closed-loop with
the controller's real-time feasibility margin rather than being an offline
manned-rotor optimization.

% =====================================================================
\section{Platform and dynamical model}\label{sec:model}

\subsection{The fully-tilting hexacopter}\label{sec:config}
The vehicle has six rotors on booms of
radius $R_r=\SI{0.70}{\meter}$ about the centre of mass. Each rotor is carried
by a nacelle that rotates independently about the body spanwise axis
$y_b$. The tilt angle $\delta_i$ is measured from the vertical:
$\delta_i=0$ in hover (thrust along $+z_b$) and $\delta_i=\pi/2$ in cruise
(thrust along $+x_b$). A fixed wing with two independent ailerons and a
V-tail with two ruddervators (mixed elevator/rudder action) provide cruise
lift and control. The six throttles, six independent tilt servos, and four
surfaces give sixteen control inputs for six DOF. The nominal parameters are
given in Table~\ref{tab:params}. Throughout the paper the dynamics are written
in an inertial $z$-up convention for clarity; the implementation uses a
body FRD convention in which the vertical rotor force is $-T\cos\delta$, which
is a sign-only difference.

\begin{table}[t]
\centering
\caption{Nominal parameters of the \SI{30}{\kilogram} fully-tilting hexacopter.}
\label{tab:params}
\small
\begin{tabular}{lll}
\toprule
Symbol & Value & Description \\
\midrule
$m$                  & \SI{30}{\kilogram}            & Maximum take-off mass \\
$W=mg$               & \SI{294.2}{\newton}           & Weight \\
$b$                  & \SI{3.50}{\meter}             & Wing span \\
$S$                  & \SI{1.26}{\meter\squared}     & Wing area \\
$AR=b^2/S$           & 9.72                          & Aspect ratio \\
$\bar c=S/b$         & \SI{0.36}{\meter}             & Mean aerodynamic chord \\
$D_p$                & \SI{0.70}{\meter}             & Propeller diameter \\
$A_d=\pi D_p^2/4$    & \SI{0.385}{\meter\squared}    & Rotor disk area (each) \\
$T_{i,\max}$         & \SI{95}{\newton}              & Per-rotor max thrust ($T/W=1.94$) \\
$T_{i,\mathrm{hov}}$ & \SI{49.0}{\newton}            & Per-rotor hover thrust ($W/6$) \\
$P_{\mathrm{ind,hov}}$ & \SI{2.12}{\kilo\watt}        & Ideal induced hover power (all rotors) \\
$V_c$                & \SI{20}{\meter\per\second}    & Test cruise speed (25 nominal) \\
$V_s$                & \SI{16.2}{\meter\per\second}  & Wing-borne stall speed ($C_{L,\max}=1.45$) \\
$V_{\max}$           & \SI{25}{\meter\per\second}    & Nominal maximum speed \\
$q_c=\tfrac12\rho V_c^2$ & \SI{245}{\pascal}         & Cruise dynamic pressure ($V_c=20$) \\
$\dot\delta_{\max}$  & \SI{60}{\degree\per\second}  & Independent tilt rate limit \\
$R_r$                & \SI{0.80}{\meter}             & Rotor boom radius \\
$I_{xx},I_{yy},I_{zz}$ & $4.27,\,6.64,\,9.58$ \si{\kilogram\meter\squared} & Body inertias \\
Tail                 & V-tail, $35^\circ$ dihedral, $S_t=0.18$ m$^2$ & Two ruddervators ($\pm25^\circ$) \\
\midrule
Actuators            & $6$ throttle $+\,6$ tilt $+\,4$ surfaces & 16 inputs, 6 DOF \\
\bottomrule
\end{tabular}
\end{table}

\subsection{Rigid-body dynamics}\label{sec:rigid}
Let $p\in\mathbb{R}^3$ and $v$ denote position and velocity in the inertial
frame, $\eta=(\phi,\theta,\psi)$ the Euler angles, and $\omega^b$ the body
rates. With $R_b^I$ the body-to-inertial rotation matrix,
\begin{equation}
\dot p = v,\qquad
m\dot v = mg_0 + R_b^I\!\left(F_r+F_a\right),\qquad
\dot\eta = H(\eta)\,\omega^b,\qquad
J\dot\omega^b = -\omega^b\times J\omega^b + \tau_r+\tau_a ,
\label{eq:rigid}
\end{equation}
where $g_0=[0,0,-g]^\top$, $F_r,\tau_r$ are the rotor forces and moments,
$F_a,\tau_a$ the aerodynamic forces and moments, and $J=\mathrm{diag}(I_{xx},I_{yy},I_{zz})$.

\subsection{Rotor and independent tilt model}\label{sec:rotor}
Rotor $i$ at body position $r_i=[x_i,y_i,z_i]^\top$ produces thrust $T_i$
along its tilted axis
\begin{equation}
n_i(\delta_i)=[\sin\delta_i,\ 0,\ \cos\delta_i]^\top .
\label{eq:naxis}
\end{equation}
The aggregate rotor force and moment are
\begin{equation}
F_r=\sum_{i=1}^{6} T_i\, n_i(\delta_i), \qquad
\tau_r=\sum_{i=1}^{6}\left[r_i\times T_i n_i(\delta_i)
   + \sigma_i \kappa T_i\, n_i(\delta_i)\right],
\label{eq:rotor}
\end{equation}
where $\sigma_i\in\{+1,-1\}$ is the rotor spin direction and $\kappa$ is the
thrust-to-reaction-torque coefficient. Throttle and tilt are rate limited,
\begin{equation}
0\le T_i\le T_{i,\max},\quad |\dot T_i|\le \dot T_{\max},\quad
0\le\delta_i\le\pi/2,\quad |\dot\delta_i|\le\dot\delta_{\max}.
\label{eq:actlim}
\end{equation}
Independence of the six $\delta_i$ is the key difference from a
synchronized-tilt or front-tilt machine: differential tilt generates direct
pitch and lift-redistribution moments, and a failed rotor can be isolated by
re-tilting the remaining healthy rotors.

\subsection{Aerodynamic and slipstream model}\label{sec:aero}
The wing and tail are modelled with stability derivatives,
\begin{equation}
F_a=\tfrac12\rho V^2 S
 \begin{bmatrix} -C_D(\alpha,\delta)\\ C_Y(\beta,\delta_r)\\ -C_L(\alpha,\delta)\end{bmatrix},
\qquad
\tau_a=\tfrac12\rho V^2 S
 \begin{bmatrix} b\, C_l(\beta,p,\delta_a)\\ \bar c\, C_m(\alpha,q,\delta_e)\\ b\, C_n(\beta,r,\delta_r)\end{bmatrix},
\label{eq:aero}
\end{equation}
with $C_L=C_{L_\alpha}\alpha+C_{L_{\delta_e}}\delta_e$ and a quadratic drag
polar $C_D=C_{D_0}+k C_L^2$. Because the rotors blow the wing during
conversion, an effective slipstream velocity and an augmented
$C_{L,\max}(\delta)$ are used, following the propeller-wing interaction
modelling common in tilt-wing studies~\cite{rohr2021nmpc,mancinelli2025unified}.
Aerodynamic control authority scales with $q=\tfrac12\rho V^2$ and therefore
vanishes in hover, where the rotors provide all control.

\subsection{Two-model strategy}\label{sec:twomodel}
Two models are deliberately maintained. A low-order model
(Eqs.~\eqref{eq:rigid}--\eqref{eq:aero}) with coarse stability derivatives is
used by the controller and the corridor OCP, while the ArduPilot SITL flight
dynamics model (higher fidelity, with motor and servo dynamics) acts as the
plant. The mismatch between the two is the principal robustness test in
Section~\ref{sec:campaign}; it is not hidden by using the same model for
control and simulation.

% =====================================================================
\section{Time-varying attainable-wrench-set control allocation}\label{sec:aws}

\subsection{Control effectiveness matrix}\label{sec:Bmatrix}
Collect the virtual wrench and the actuators as
\begin{equation}
w=\begin{bmatrix}F\\ \tau\end{bmatrix}\in\mathbb{R}^{6},\qquad
u=\begin{bmatrix}T_1\!\cdots T_6&\delta_1\!\cdots\delta_6&\delta_{aL}&\delta_{aR}&\delta_{rL}&\delta_{rR}\end{bmatrix}^\top\in\mathbb{R}^{16}.
\label{eq:wu}
\end{equation}
Because the rotor force contains products $T_i\sin\delta_i$ and
$T_i\cos\delta_i$, the mapping is nonlinear. Linearizing about the operating
point $(T_0,\delta_0,V)$ gives
\begin{equation}
w \approx w_0 + B(\delta_0,V)\,\Delta u,\qquad
B(\delta_0,V)=\begin{bmatrix}
B_{FT}(\delta_0) & B_{F\delta}(\delta_0) & 0\\
B_{\tau T}(\delta_0) & B_{\tau\delta}(\delta_0) & B_{\tau s}(V)
\end{bmatrix},
\label{eq:B}
\end{equation}
where the surface block $B_{\tau s}(V)\propto q(V)$ grows with airspeed and the
rotor blocks rotate continuously with $\delta_0$. At hover $B_{\tau s}\approx0$
and the rotors act in all six channels; in cruise $B_{\tau s}$ dominates the
moment channels while the rotors act mainly along $x_b$. The matrix is
re-evaluated at every control step.

\subsection{Attainable wrench set and inscribed-sphere metric}\label{sec:awsmetric}
The attainable wrench set at the current operating point is
\begin{equation}
\mathcal{M}(\delta,V)=\bigl\{w:\ w=w_0+B(\delta,V)\Delta u,\
\underline u\le u\le\bar u,\ |\dot u|\le\dot u_{\max}\bigr\},
\label{eq:AWS}
\end{equation}
a bounded six-dimensional polytope. Following the passive fault-tolerance
construction of~\cite{yang2026biaxial}, its size is summarized by the radius
of the largest wrench ball centered at the nominal $w_0$ that it contains,
\begin{equation}
\rho(\delta,V)=\max_{r\ge0}\ r\quad\text{s.t.}\quad
\mathcal{B}_r(w_0)\subseteq\mathcal{M}(\delta,V).
\label{eq:rho}
\end{equation}
A large $\rho$ means ample authority in every direction; $\rho\to0$ signals
incipient saturation; and a commanded wrench with
$\|w_{\mathrm{cmd}}-w_0\|_2>\rho$ is infeasible. In this work the normalized
minimum singular value of the effectiveness matrix is used as a scalar proxy
for $\rho$ and is evaluated along the corridor to characterize authority
(Fig.~\ref{fig:aws}); the constrained QP of Eq.~\eqref{eq:alloc} enforces the
exact actuator limits online, so the explicit inscribed-sphere LP is not on
the critical real-time path.

\subsection{Online weighted allocation}\label{sec:alloc}
At each fast-loop step the MPC produces $w_{\mathrm{cmd}}$. The actuator
command solves a weighted, prioritized least-squares
problem~\cite{johansen2013survey,blaha2023survey,tohidi2020adaptive},
\begin{equation}
\min_{\Delta u}\ \bigl\|\Omega_w(B\Delta u-\Delta w_{\mathrm{cmd}})\bigr\|_2^2
 +\bigl\|\Omega_u(\Delta u-\Delta u_d)\bigr\|_2^2
\quad\text{s.t.}\quad \underline u\le u\le\bar u,\ |\dot u|\le\dot u_{\max},
\label{eq:alloc}
\end{equation}
with fixed weights $\Omega_w,\Omega_u$ chosen so that the six independent tilt
angles are used preferentially for moment generation at low speed and the
surfaces at high speed, consistently with the time-varying $B(\delta,V)$.
Because the limits are hard constraints, the active-set solver redistributes
the wrench to the healthy channels whenever an actuator reaches its bound;
this is the mechanism that provides online load redistribution, rather than an
explicit adaptive re-weighting. The achieved wrench and the unmet residual are
returned at every step and used to freeze the error integral under
saturation.

\subsection{Feasibility monitoring and wrench hedging}\label{sec:hedge}
The primary online feasibility mechanism is the constrained QP of
Eq.~\eqref{eq:alloc}: actuator position and rate limits are hard constraints,
so the allocated wrench is always attainable and the residual
$w_{\mathrm{cmd}}-w_{\mathrm{ach}}$ is explicitly returned rather than silently
clipped. In addition, the inscribed-sphere margin $\rho$ is evaluated along the
realised trajectory as a scalar feasibility monitor. If the commanded wrench
ball were to exceed $\rho$, a wrench-hedging projection onto the boundary of
$\mathcal{M}$ and a request to slow the outer reference would be issued. This
projection and the associated outer re-planning are part of the architecture
but, as reported in Section~\ref{sec:campaign}, they were never triggered in
the tested envelope because the corridor reference together with the
constrained QP kept $\rho$ strictly positive; the realised minimum margin is
reported rather than asserted. This design replaces the abrupt saturation and
integrator wind-up of a fixed allocator and defines the signal that would
close the two-layer loop of Section~\ref{sec:twolayer} under a more demanding
command.

% =====================================================================
\section{Corridor-aware optimal transition scheduling}\label{sec:corridor}

\subsection{Conversion corridor}\label{sec:corridorderiv}
The lower (stall) boundary follows from the vertical force balance. With the
rotors contributing a vertical component $\sum_i T_i\cos\delta_i$ and the wing
providing $L=\tfrac12\rho V^2 S C_L$, the minimum sustainable speed is
\begin{equation}
V_{\min}(\delta)=
\sqrt{\frac{2\,\bigl(mg-\sum_i T_{i,\max}\cos\delta\bigr)}
{\rho\,S\,C_{L,\max}(\delta)}},
\qquad V_{\min}=0\ \text{if }mg\le\sum_i T_{i,\max}\cos\delta .
\label{eq:Vmin}
\end{equation}
For the \SI{30}{\kilogram} vehicle, $\sum_i T_{i,\max}=6\times95=\SI{570}{\newton}$
in total. The corridor restricts incidence to a conservative
$\alpha_{\mathrm{stall}}=13^\circ$, which gives
$C_{L,\mathrm{stall}}=C_{L_0}+C_{L_\alpha}\alpha_{\mathrm{stall}}=1.29$;
the resulting boundary at $\delta=90^\circ$ is $17.2$~m/s, slightly above the
$C_{L,\max}=1.45$ wing-borne stall speed $V_s=16.2$~m/s listed in
Table~\ref{tab:params}, because the corridor does not use the full lift
coefficient. The upper boundary is set by the
available power and the normal-load limit,
\begin{equation}
P_{\mathrm{avail}}\ge P_{\mathrm{ind}}+P_{\mathrm{prof}}+P_{\mathrm{par}}
 +V\!\sum_i T_i\sin\delta_i,\qquad |n|\le n_{\max},
\label{eq:Vmax}
\end{equation}
which defines $V_{\max}(\delta)$. The corridor is
$\mathcal{C}=\{(\delta,V):V_{\min}(\delta)\le V\le V_{\max}(\delta)\}$; its
nominal lower boundary is listed in Table~\ref{tab:corridor} from
Eq.~\eqref{eq:Vmin} with the conservative, no-slipstream stall incidence
$\alpha_{\mathrm{stall}}=13^\circ$, i.e. $C_{L,\mathrm{stall}}=1.29$ rather
than the full $C_{L,\max}=1.45$ (a soft-stall model
$C_L=C_{L,\max}\tanh((C_{L,0}+C_{L,\alpha}\alpha)/C_{L,\max})$ is used
identically in the plant, the prediction model, and the corridor).

\begin{table}[t]
\centering
\caption{Nominal lower conversion-corridor boundary for the \SI{30}{\kilogram}
vehicle from Eq.~\eqref{eq:Vmin}, using the conservative corridor stall
incidence $\alpha_{\mathrm{stall}}=13^\circ$ ($C_{L,\mathrm{stall}}=1.29$, no
slipstream augmentation). The $\delta=90^\circ$ value ($17.2$~m/s) is therefore
slightly above the $C_{L,\max}=1.45$ wing-borne stall speed $V_s=16.2$~m/s.}
\label{tab:corridor}
\small
\begin{tabular}{lccccccc}
\toprule
Tilt $\delta$ (deg) & 0 & 15 & 30 & 45 & 60 & 75 & 90 \\
\midrule
$V_{\min}$ (m/s) & 0 & 0 & 0 & 0 & 3.0 & 12.1 & 17.2 \\
\bottomrule
\end{tabular}
\end{table}

\subsection{Forward transition reference}\label{sec:ocp}
On a level, slow conversion the full dynamic optimal-control problem with the
power/load/tilt-rate objective collapses to a quasi-steady minimum-thrust
trim traversed by a jerk-limited speed law; we therefore use this explicit,
deterministic form rather than a free-time collocation solver (a CasADi/IPOPT
formulation~\cite{andersson2019casadi} was implemented and discarded because it
is singular at
$V=0$ and produces bang--bang tilt activity). A smoothstep speed S-curve
$V(s)=V_0+(V_f-V_0)(3s^2-2s^3)$, $s=t/T_f$, with zero end jerk, prescribes
the corridor traversal in a fixed conversion time $T_f$. At each node the
\emph{minimum-thrust level trim} is solved from the exact longitudinal
force-and-moment equilibrium
\begin{equation}
\sum_i T_i\sin\tau = D,\qquad
\sum_i T_i\cos\tau + L = mg,\qquad
M_y^{\mathrm{rot}}+M_y^{\mathrm{wing}}+M_y^{\mathrm{tail}}=0,
\label{eq:trim}
\end{equation}
with $\tau=\beta-\theta$ the thrust-line angle from vertical, including the
thrust-line moment and the V-tail authority, and with the wing loaded up to
exactly what is needed (bounded by $C_{L,\max}$); the rotors supply the
residual lift and the forward force. The node solution
$(\beta,T,\theta,\alpha)$ and the body-frame feed-forward wrench
$w_{\mathrm{ff}}$ and aerodynamic-surface force $F_{\mathrm{surf}}$ are
tabulated against airspeed. The stall boundary $V_{\min}(\beta)$ is enforced
at every node, so the reference is feasible by construction.

\subsection{Backward transition}\label{sec:backward}
Backward transition is treated symmetrically but is the more sensitive
manoeuvre, because premature wing unloading at intermediate speed produces a
pitch-up and a sink~\cite{may2025backward}. Indexing the trim on the
\emph{measured} airspeed was found to be unstable in our plant; instead the
backward leg uses a monotonic, time-scheduled rotor-borne incidence profile
(cruise incidence easing from $11^\circ$ to $0$ as the nacelles tilt from
$63^\circ$ back to $0^\circ$), which keeps the rotors dominant as speed falls
and removes the near-stall pitch-up. Braking is obtained by reducing forward
thrust (drag and thrust removal); an aft-thrust feed-forward is clamped to
zero, since it drives the nacelles aft at low speed and is destabilising.

\subsection{Trajectory library and online re-planning}\label{sec:library}
The trim map and the two traversals are computed once offline over a grid of
mass, cruise speed, and wind conditions, producing a library of references
$(x^\star,V^\star,h^\star,\theta^\star)$ together with the feed-forward
wrench and surface force. The online layer indexes the forward leg by the
reference speed and holds the cruise node's trim ($\theta=11^\circ$,
$\beta=63^\circ$) at the entry to the backward leg so that the reference is
$C^0$-continuous across the cruise/backward boundary; re-planning inside the
corridor is requested when the feasibility feedback of
Section~\ref{sec:hedge} falls below its threshold.

% =====================================================================
\section{Unified MPC and two-layer integration}\label{sec:nmpc}

\subsection{Unified error-space LTV-MPC formulation}\label{sec:nmpcform}
A single controller covers hover, transition, and cruise without mode
switching. The nonlinear plant is linearized, at each control step, in error
coordinates about the speed-dependent nonlinear trim reference of
Section~\ref{sec:corridor}; the resulting linear time-varying (LTV) problem is
then discretized and solved as a quadratic program. This is an error-space
LTV-MPC about a nonlinear reference rather than a fully nonlinear programming
NMPC; we retain the shorthand MPC throughout. Over horizon $N$ and period
$\Delta t$ it minimizes
\begin{equation}
\min_{\{u_k\}}\ \sum_{k=0}^{N-1}
\!\left[\|x_k-x_k^\star\|_{Q}^{2}+\|u_k-u_k^\star\|_{R}^{2}\right]
+\|x_N-x_N^\star\|_{Q_f}^{2}
\label{eq:nmpc}
\end{equation}
subject to the discrete prediction model, the actuator
limits~\eqref{eq:actlim}, and a feasibility constraint
$\rho(\delta_k,V_k)\ge\rho_{\min}$. The reference $x_k^\star,u_k^\star$ is
provided by the corridor trim map of Section~\ref{sec:corridor}. The
12-state NED model is linearised in error coordinates about the
speed-dependent reference and discretised exactly by zero-order hold with a
matrix exponential (\texttt{scipy.linalg.expm}); the horizon is
$N=20$ at $\Delta t=0.03$~s ($\approx 33$~Hz). The resulting equality
constraints are condensed into a small dense quadratic program in the
five virtual-wrench increments, which is solved by the active-set
\texttt{quadprog} solver (an embedded condensed QP such as
acados~\cite{verschueren2022acados} would be the natural flight-hardware
target); on infeasibility the attitude soft constraints are
relaxed first, and a zero-wrench fallback is used only as a last resort.
The virtual wrench is then passed to the time-varying AWS allocator of
Section~\ref{sec:aws}. A bounded, model-free position-error integral
augments the wrench to reject constant plant mismatch (Section~\ref{sec:campaign}).

\subsection{Model-free position-error integral for constant mismatch}\label{sec:integral}
An acceleration-residual disturbance observer was first implemented on the
internal prediction model, but it proved unsafe during the dynamic backward
conversion: because the internal model omits the realised V-tail download and
motor dynamics, the observer mistakes the time-varying, predictable
model--plant aerodynamic difference for a disturbance and cancels the trim
feed-forward, reproducing a stall/crash at intermediate speed. We instead use
a bounded, model-free integral of the \emph{position} error
$e_p=p^\star-p$ in the NED frame,
\begin{equation}
\mathcal I_k = \operatorname{sat}\!\bigl(\mathcal I_{k-1}+K_i\,e_p\,\Delta t\bigr),
\qquad
w_k = w_k^{\mathrm{MPC}}+R_b^\top\mathcal I_k ,
\label{eq:integral}
\end{equation}
which is rotated into the body frame and added to the virtual wrench before
allocation. The vertical channel integrates only during cruise, backward
conversion, and terminal hover ($K_{i,z}=4$, $\pm70$~N) so that a thrust or
mass deficit is corrected as the wing unloads; the horizontal channels
integrate only in terminal hover ($K_{i,h}=7$, $\pm80$~N) to trim a steady
crosswind, and are frozen in climb and conversion to avoid wind-up against a
moving reference. The integral is frozen whenever the allocator reports
saturation. This adds no model knowledge and is the mechanism that recovers
the constant-mismatch robustness cases of Table~\ref{tab:robust}.

\subsection{Two-layer architecture and feasibility coupling}\label{sec:twolayer}
The two layers are separated by time scale and coupled through feasibility:
\begin{itemize}[leftmargin=1.6em,itemsep=2pt]
  \item \emph{Outer layer (1--10\,Hz):} corridor OCP, using the nominal
  corridor and a nominal AWS, outputs the conversion reference.
  \item \emph{Inner layer (20--50\,Hz):} unified MPC plus the time-varying
  AWS allocation of Section~\ref{sec:aws}, tracking the reference against the
  real (disturbed) plant and evaluating the true $\rho$.
  \item \emph{Coupling (designed):} when $\rho<\rho_{\mathrm{th}}$ (due to a
  gust, mass error, or saturation), the inner layer is designed to hedge
  immediately and request a slower, more conservative outer reference, which
  the outer layer would re-plan inside the corridor. In the campaign reported
  here the threshold was never crossed (the realised $\rho$ stayed positive),
  so the closed-loop re-planning is a designed but not exercised path; the
  feasibility that was actually demonstrated online is provided by the
  constrained QP at every step and by the feasible-by-construction corridor
  reference. Triggering and validating the re-planning path under deliberately
  infeasible commands is left to future work.
\end{itemize}

\subsection{Single-stick velocity command}\label{sec:stick}
A uniform command mapping is used across the envelope: pitch stick to
$v_x$, roll stick to $v_y$, throttle to $v_z$, and rudder to $\dot\psi$.
Following the simplified-operation motivation of~\cite{milz2026taming} and the
reference-model idea of~\cite{bhardwaj2018integrated}, this is treated as a
design property rather than a novelty claim. The evaluated objective metric is
the number of explicit mode switches required to complete a full mission: the
stock switched architecture enters and leaves the fixed-wing controller
several times, whereas the unified controller switches zero times
(Table~\ref{tab:main}). A quantitative stick-to-velocity transfer-function
identification across the envelope is a planned extension and is not claimed
here.

\subsection{Real-time implementation}\label{sec:realtime}
The inner MPC is warm-started from the previous solution and the dense QP
uses a fixed iteration cap; the solve time, deadline misses, and the fallback
path (infeasible $\rightarrow$ relaxed attitude constraints $\rightarrow$
zero-wrench hold) are logged at every step. Because the condensed problem has
only $5N$ decision variables, the mean/P99/worst-case solve time over the
whole campaign is $0.5$/$2.1$/$5.5$~ms on a single desktop core, against a
$30$~ms control period (Section~\ref{sec:campaign}), i.e.\ the controller is
two orders of magnitude away from its deadline. The outer trim map is
computed offline and is a table lookup online, so it has negligible runtime.

% =====================================================================
\section{ArduPilot SITL implementation}\label{sec:sitl}

\subsection{Stock tilt-rotor architecture and its limitation}\label{sec:stock}
ArduPilot models tilt-rotors as a special QuadPlane: a six-rotor frame is
selected with \texttt{Q\_FRAME\_CLASS=2}, the tilting motors are declared with
\texttt{Q\_TILT\_MASK}, and the tilt type is set to continuous
(\texttt{Q\_TILT\_TYPE=0}) with \texttt{Q\_TILT\_MAX},
\texttt{Q\_TILT\_RATE\_UP}, and \texttt{Q\_TILT\_RATE\_DN} defining the
schedule~\cite{ardupilot2024tilt,ardupilot2024transitions}. In this mode the
multirotor and fixed-wing controllers are switched, and the tilt servos are
commanded as a single synchronized logical tilt following a fixed rate. This
stock behaviour is the baseline controller \emph{and} it exposes a genuine
limitation: independent per-rotor tilt angles are not represented by the
stock continuous-tilt scheduler.

\subsection{Companion-layer integration}\label{sec:companion}
To exploit six independent tilts, the proposed controller is implemented as a
Python companion that closes the loop around a high-rate nonlinear plant
that reproduces the ArduPilot SITL sensor/actuator interfaces: the plant
advances at 400~Hz with a 4th-order Runge--Kutta integrator (two substeps),
while the MPC and allocator run at 33~Hz. All quantitative results in
Section~\ref{sec:campaign} are generated by this companion plant. The
control-effectiveness, constraints, and active-set logic are written in C++
as a reusable library (\texttt{AP\_TiltHexa}) and called from the companion
through a small C ABI (\texttt{ctypes}); the same C++ core is the candidate
for in-firmware deployment. The plant exposes the 16 actuator channels (six
throttles, six independent tilt servos, and the four V-tail/aileron
surfaces) and logs the full truth state, but the controller receives only
the sensed state; truth columns are used exclusively for post-processing.
The 16-channel V-tail configuration is declared explicitly. A small,
publicly diffed firmware hook exposes the six tilt servos and the surfaces
for the native \texttt{arduplane} SITL target; the benchmark can be
reproduced from the companion alone without rebuilding the firmware, while
porting the closed loop onto the native \texttt{arduplane} binary is an
integration step that is not part of the reported numbers.

\subsection{Baseline and fairness}\label{sec:baseline}
The comparison baseline is an incremental-nonlinear-dynamic-inversion
(INDI) attitude/rate controller with the earlier constrained
weighted-least-squares / quadratic-programming allocator (INDI--WLS), which is
the conventional engineering solution for over-actuated convertible aircraft.
Both controllers receive the identical mission, disturbance, and plant, and
neither reads the truth state; the baseline parameters are fixed before the
campaign. A second ablation, \emph{proposed without the corridor}, uses the
same MPC but a fixed-rate tilt schedule, isolating the contribution of the
corridor reference.

% =====================================================================
\section{Simulation campaign}\label{sec:campaign}

\subsection{Scenario matrix}\label{sec:scenarios}
The same scenarios are flown by the baseline and the proposed controller,
against the identical plant, disturbance, and mission. The primary scenario is
a complete hover--cruise--hover mission at nominal mass, which contains both
the forward and backward conversion legs and therefore serves as the nominal
forward- and backward-transition case; the remaining scenarios are robustness
variations: (S4) the nominal mission with a gentle coordinated turn
($5^\circ$/s yaw rate, $\approx 10^\circ$ bank); (S5) a 5\,m/s discrete gust;
(S6) a 4\,m/s steady crosswind; (S7) mass $+15\%$; (S8) a forward
centre-of-gravity offset of 0.025\,m; and (S9) a control-model mismatch sweep:
(S9a) thrust $-10\%$, (S9b) control-surface effectiveness $-20\%$, and (S9c)
inertia $+20\%$. A dedicated hover and attitude-step case is not part of the
reported campaign and is left to future work. A 20-run Monte-Carlo campaign
randomises the same parameters jointly (seeds 1000--1019). A run is accepted
only on nine plant-truth checks (altitude and cruise reached, attitude below
$60^\circ$, no post-climb altitude loss, inside the corridor, non-negative
allocation residual, symmetric conversions, monotonic phasing with zero mode
switch, and terminal hover judged by ground speed).

\subsection{Metrics}\label{sec:metrics}
Performance is measured by: transition time $t_f$; maximum altitude loss
$\Delta h_{\max}$; peak and RMS attitude error; total control energy
$\int u^\top R_u u\,\mathrm dt$; path tracking RMS; number of mode switches;
and the AWS margin minimum $\rho_{\min}$ along the transition. Real-time
performance is measured by the MPC solve-time mean, 99th percentile, and
worst case, and by the deadline-miss rate.

\subsection{Results protocol}\label{sec:results}
All quantitative results are recomputed from the plant-truth logs by the
evaluation script; the controller never reads the truth. The corridor and
trim schedules are shown in Figs.~\ref{fig:corridor}--\ref{fig:trim}, the
nominal full-profile time history in Fig.~\ref{fig:profile}, the corridor
ablation in Fig.~\ref{fig:ablation}, the rotor effectiveness metric in
Fig.~\ref{fig:aws}, and the tracking and real-time robustness summaries in
Figs.~\ref{fig:robust}--\ref{fig:rt}. Comparative numbers are reported in
Tables~\ref{tab:main}--\ref{tab:rt}.

\begin{figure}[t]
\centering\includegraphics[width=0.62\columnwidth]{figures/fig_corridor.pdf}
\caption{Conversion corridor: stall boundary $V_{\min}(\beta)$, feasible
region, and the forward/backward trim schedules.}
\label{fig:corridor}
\end{figure}

\begin{figure}[t]
\centering\includegraphics[width=\columnwidth]{figures/fig_trim_schedule.pdf}
\caption{Forward/backward minimum-thrust trim: nacelle tilt, per-rotor
thrust, and body pitch versus airspeed.}
\label{fig:trim}
\end{figure}

\begin{figure*}[t]
\centering\includegraphics[width=0.92\textwidth]{figures/fig_profile.pdf}
\caption{Nominal hover--cruise--hover mission (truth): altitude, airspeed,
attitude, mean nacelle tilt, mean rotor thrust, and ruddervator deflection.
Shaded bands mark the six mission phases.}
\label{fig:profile}
\end{figure*}

\begin{figure}[t]
\centering\includegraphics[width=\columnwidth]{figures/fig_ablation_corridor.pdf}
\caption{Corridor ablation: airspeed (left) and altitude (right) for the
corridor-aware MPC and the fixed-schedule ablation. The ablation overshoots
cruise speed and loses 7\,m on the backward leg.}
\label{fig:ablation}
\end{figure}

\begin{figure}[t]
\centering\includegraphics[width=0.62\columnwidth]{figures/fig_aws.pdf}
\caption{Normalised minimum singular value of the rotor effectiveness block
versus airspeed (the AWS inscribed-sphere metric).}
\label{fig:aws}
\end{figure}

\begin{figure}[t]
\centering\includegraphics[width=\columnwidth]{figures/fig_robustness.pdf}
\caption{Altitude and speed tracking RMSE across the fixed robustness
campaign; all scenarios pass the nine truth checks.}
\label{fig:robust}
\end{figure}

\begin{figure}[t]
\centering\includegraphics[width=\columnwidth]{figures/fig_mc.pdf}
\caption{Monte-Carlo campaign (20 seeds, joint parameter perturbation). Left:
altitude and speed tracking RMSE distributions; right: per-seed P99 and worst
MPC solve time. All 20 runs pass the nine truth checks.}
\label{fig:mc}
\end{figure}

\begin{figure}[t]
\centering\includegraphics[width=\columnwidth]{figures/fig_solve_time.pdf}
\caption{MPC solve-time mean, P99, and worst case across the fixed campaign,
all far below the 30\,ms control period.}
\label{fig:rt}
\end{figure}

\begin{table}[t]
\centering
\caption{Main transition results from the plant-truth logs. Energy is the
rotor induced energy $E=\int\sum_i T_i^{3/2}/\sqrt{2\rho A}\,\mathrm dt$ over
each conversion (momentum theory, rotor diameter $0.70$~m). The fixed-schedule
ablation overshoots the cruise speed (24.5 vs 20~m/s) and loses $7$~m on the
backward leg.}
\label{tab:main}
\small
\begin{tabular}{lcccccc}
\toprule
 & \multicolumn{3}{c}{Forward transition} & \multicolumn{3}{c}{Backward transition} \\
\cmidrule(lr){2-4}\cmidrule(lr){5-7}
Controller & $t_f$ (s) & $\Delta h_{\max}$ (m) & Energy (kJ) & $t_f$ (s) & $\Delta h_{\max}$ (m) & Energy (kJ) \\
\midrule
Proposed (full)        & 14.0 & 1.59 & 21.0 & 20.1 & 0.98 & 26.4 \\
Proposed, no corridor  & 14.0 & 2.67 & 19.9 & 20.1 & 7.00 & 45.6 \\
INDI--WLS baseline     & 10.0 & 0.08 & 18.6 & 15.1 & 0.05 & 27.0 \\
\bottomrule
\end{tabular}
\end{table}

\begin{table}[t]
\centering
\caption{Robustness campaign (proposed controller, plant-only mismatch); every
scenario passes all nine truth checks. Terminal ground speed is used for the
hover criterion (in a steady crosswind a correct hover has airspeed equal to
the wind speed). All runs use zero explicit mode switches.}
\label{tab:robust}
\small
\begin{tabular}{lcccccc}
\toprule
Scenario & $h_f$ (m) & $V_f$ (m/s) & $h$ RMSE (m) & $V$ RMSE (m/s) & $\max|\phi|$ (deg) & $\max|\theta|$ (deg) \\
\midrule
Nominal            & 5.07 & 0.18 & 0.84 & 0.48 & 0.0 & 11.2 \\
S4 coordinated turn & 5.12 & 0.57 & 0.81 & 0.49 & 14.4 & 12.2 \\
S5 5\,m/s gust      & 5.16 & 0.66 & 0.91 & 0.59 & 0.0 & 17.6 \\
S6 4\,m/s crosswind & 5.26 & 0.60 & 1.08 & 0.43 & 3.6 & 11.1 \\
S7 mass $+15\%$     & 5.19 & 0.69 & 1.71 & 0.43 & 0.0 & 12.0 \\
S8 CG forward 0.025\,m & 5.12 & 0.67 & 0.81 & 0.53 & 0.0 & 11.1 \\
S9a thrust $-10\%$  & 5.20 & 0.59 & 1.35 & 0.38 & 0.0 & 11.0 \\
S9b surface $-20\%$ & 5.14 & 0.58 & 0.82 & 0.49 & 0.0 & 11.0 \\
S9c inertia $+20\%$ & 5.13 & 0.56 & 0.82 & 0.48 & 0.0 & 11.2 \\
\bottomrule
\end{tabular}
\end{table}

\begin{table}[t]
\centering
\caption{Real-time performance on one desktop core (control period 30\,ms).
The inner MPC is the condensed QP ($N=20$, five virtual-wrench channels);
the allocator is the C++ active-set QP called through a C ABI; the trim map is
computed offline and is a table lookup online. Ranges are over the fixed
robustness campaign.}
\label{tab:rt}
\small
\begin{tabular}{lcccc}
\toprule
Stage & Mean & P99 & Worst & Deadline-miss rate \\
\midrule
Inner MPC (33\,Hz) & 0.41--0.56\,ms & 1.5--2.4\,ms & 5.5\,ms$^{a}$ & 0\%$^{a}$ \\
AWS allocator (C++) & 49--74\,$\mu$s & --- & --- & 0\% \\
Outer trim map & \multicolumn{4}{c}{offline; online table lookup} \\
\bottomrule
\end{tabular}\\[2pt]
\footnotesize $^{a}$Over the fixed campaign. One isolated solve of 37.6\,ms
(one of $\approx$40\,000 Monte-Carlo solves) exceeded the 30\,ms period, a
single desktop scheduling hiccup; the P99 remained below 2.4\,ms.
\end{table}

The campaign confirms the following. The corridor-aware controller completes
the full mission with zero explicit mode switches and peak attitude below
$18^\circ$ (limit $60^\circ$). The online feasibility indicators support the
framing of Section~\ref{sec:hedge}: the smallest corridor margin on the
realised trajectory is $0.72$~m/s, the maximum unallocated wrench residual is
$0.038$ in mixed units (mean $10^{-3}$), and the thrust, tilt, and surface
saturation fractions are all zero, so the wrench-hedging and outer
re-planning path was correctly never entered. Removing the corridor (fixed tilt schedule)
leaves the aircraft flying but it overshoots cruise speed to 24.5\,m/s,
increases the backward-leg altitude deviation from $0.98$ to $7.0$~m, and
nearly doubles the backward conversion energy (45.6 vs 26.4\,kJ), so the
corridor is doing real work rather than merely reshaping the reference. The
bounded position integral rejects constant mass, thrust, and wind mismatch
without the instability of an acceleration-residual observer. The inner QP
runs at a mean of about $0.5$~ms (P99 below $2.4$~ms), and the C++ allocator
at $49$--$74\,\mu$s, both comfortably real-time on a desktop core. The
envelope has clear limits, reported honestly: a $15^\circ$/s, $28^\circ$-bank
coordinated turn is not trackable (the aircraft sideslips and descends), so a
gentle $5^\circ$/s turn is used; a forward CG offset of 0.05\,m or more
produces a takeoff pitch runaway, and the robust envelope is 0.025\,m.

The 20-run Monte-Carlo campaign (joint mass, inertia, thrust, surface, CG,
wind, delay, and sensor-noise perturbation) passes all nine checks on every
seed: altitude RMSE mean 1.06\,m (worst 1.81\,m), speed RMSE mean 0.54\,m/s
(worst 0.77\,m/s), with the solve-time P99 below 2.4\,ms on every seed
(Fig.~\ref{fig:mc}).

The comparison against the INDI--WLS baseline is reported honestly rather
than as a foregone victory. On the same plant, with both controllers denied
the truth state, a well-tuned INDI--WLS is a strong performer: it adopts a
predominantly wing-borne conversion (nacelles reaching only about
$23^\circ$ at 20\,m/s, pitch about $5^\circ$), which converts faster
(10.0/15.1\,s vs 14.0/20.1\,s) and holds altitude about an order of
magnitude tighter ($\Delta h_{\max}$ 0.08/0.05\,m vs 1.59/0.98\,m), with
comparable conversion energy and speed RMSE (Table~\ref{tab:main}); it also
passed every S5--S9 disturbance it was tested on (gust, crosswind,
$+15\%$ mass, CG offset, thrust $-10\%$, surface $-20\%$), with a
transition-altitude RMSE of $0.04$--$0.05$~m and no altitude loss, and
remained corridor-safe even when the commanded acceleration was more than
doubled. We do not claim the proposed controller dominates a carefully tuned
INDI--WLS on tracking accuracy in the tested envelope. The two controllers differ in what they guarantee. The
INDI--WLS conversion relies on a hand-scheduled speed ramp and a
discrete hover/transition/cruise state machine, and its success is empirical:
there is no pre-maneuver certificate that the commanded tilt stays above the
stall boundary. The proposed controller, by contrast, runs a single
optimization over the entire profile (zero explicit mode switches) and
constructs the conversion reference from the offline corridor OCP, so
feasibility against the stall boundary and actuator limits is built into the
reference and enforced by the constrained QP at every step. The decisive
controlled experiment is the corridor ablation: replacing the feasibility-aware
reference with a fixed tilt schedule causes the same MPC to overshoot cruise
speed and lose $7$~m on the backward leg, which localises the benefit to the
feasibility layer rather than to optimisation per se. The tighter altitude
hold of the wing-borne baseline also indicates a concrete tuning opportunity
for the proposed scheme, namely shifting the corridor trim toward a more
wing-borne distribution at intermediate speeds, which is left as future work.

% =====================================================================
\section{Discussion and limitations}\label{sec:discussion}
The scope is deliberately simulation-only, and this must be stated plainly
against the real-flight evidence of~\cite{bauersfeld2021mpc,allenspach2021nmpc,mancinelli2025unified}.
SITL cannot reproduce several effects that dominate real convertible flight:
high-fidelity turbulent inflow and unsteady propeller-wing interaction,
servo and motor transport delays and backlash, sensor noise and estimator
drift, structural flexibility, and ground effect. The two-model strategy of
Section~\ref{sec:twomodel} and the mismatch sweep (S9) probe robustness, but
they are not equivalent to flight. The independent-tilt mixer also requires a
firmware change in mode~(b); the non-invasive mode~(a) is reproducible
without rebuilding ArduPilot but does not exercise full per-rotor tilting.
Accordingly, the contribution is framed as an open, reproducible benchmark
and a feasibility-aware allocation/planning method, not as a flight-proven
controller. Hardware-in-the-loop validation, servo-delay modelling, and
outdoor flight are the immediate next steps, followed by extension to rotor
faults, where the independent tilts and the AWS metric offer a direct
fault-tolerance path~\cite{yang2026biaxial,santos2021fast,giribet2016tilted}.

% =====================================================================
\section{Conclusions}\label{sec:conclusion}
A full-envelope control architecture was presented for a \SI{30}{\kilogram}
fully-tilting hexacopter with six independently tilting rotors. It combines
(i) a time-varying control-allocation layer built on the attainable wrench set
and its inscribed-sphere margin, with a constrained QP enforcing actuator
limits at every step, (ii) a corridor-constrained trim optimization that
supplies a feasible-by-construction conversion reference tracked by a single
unified MPC with zero mode switches, and (iii) an open SITL companion benchmark
in the ArduPilot framework with a standardized comparison against an INDI--WLS
engineering baseline. The paper deliberately does not claim unified control as
new, nor tracking superiority over the tuned INDI--WLS baseline, which converts
faster and holds altitude more tightly in nominal conditions. The stated
contributions are the full-envelope, time-varying AWS allocation for the
over-actuated fully-tilting hexacopter, the corridor-feasible reference together
with a specified feasibility-coupling path, and a reproducible open benchmark.
The conversion-corridor boundary was derived analytically for the stated
vehicle, and the corridor ablation localizes the stabilizing benefit to the
feasibility layer rather than to optimization per se. Future work will trigger
and validate the online re-planning path under deliberately infeasible
commands, shift the corridor trim toward a more wing-borne distribution, port
the closed loop to the native ArduPilot SITL target, and proceed to
hardware-in-the-loop, flight tests, and rotor-fault scenarios.

% =====================================================================
\section*{CRediT authorship contribution statement}
Author One: Conceptualization, Methodology, Software, Writing -- original draft.
Author Two: Validation, Visualization, Writing -- review \& editing.
Corresponding Author: Supervision, Funding acquisition, Writing -- review \& editing.

\section*{Declaration of competing interest}
The authors declare that they have no known competing financial interests or
personal relationships that could have appeared to influence the work reported
in this paper.

\section*{Data availability}
The ArduPilot SITL parameters, scenario files, controller source, baseline
configuration, and random seeds will be released in a public repository to
permit full reproduction of the benchmark.

\appendix
\section{Open-source artifacts and reproducibility}\label{app:artifacts}
The released artifact contains:
\begin{itemize}[leftmargin=1.6em,itemsep=2pt]
  \item ArduPilot parameter files (stock baseline and proposed configuration),
  including \texttt{Q\_FRAME\_CLASS}, \texttt{Q\_TILT\_MASK}, tilt rates, and
  limits;
  \item the nonlinear SITL plant for the \SI{30}{\kilogram} fully-tilting
  hexacopter and the Python companion controller (with a C ABI to the C++
  allocation library);
  \item the corridor/trim optimization and the error-space LTV-MPC source
  (Python with \texttt{scipy} and \texttt{quadprog}), together with the C++
  AWS allocation routine and its unit tests;
  \item the scenario scripts (nominal plus S4--S9), Monte-Carlo seeds, and plotting code;
  \item a step-by-step build/run guide and, for mode~(b), the public firmware
  diff exposing the six independent tilt servos.
\end{itemize}

\bibliographystyle{elsarticle-num}
\bibliography{references}

\end{document}
